\section{Introduction}\label{sec:intro}

Multi-object tracking (MOT) finds every object of interest in a video, such as people or cars,
and gives each object one identity number (ID) that stays the same over time. It is used in
surveillance, driving, and drone systems \cite{zhu2022visdrone,wu2022uavsurvey}, and many of these
systems must give results with little delay, so the processing cost matters as much as the
accuracy.

The performance of a MOT system depends on several factors, including the quality of object
detection, the ability to associate objects across consecutive frames, and the fit between the two.
These factors are closely related, since the detections decide what the tracker can follow. Missed
detections may interrupt object trajectories, while false detections can lead to incorrect
associations and identity switches. Achieving a suitable balance between finding the objects and
keeping their identities is therefore an important challenge in MOT systems.

Most MOT systems follow the tracking-by-detection design \cite{adzemovic2026survey}. First, an
object detector, a neural network, looks at each frame and returns candidate boxes. Each box has a
confidence score between 0 and 1 that says how sure the detector is that the box holds an object.
Then a tracker links the boxes of consecutive frames into tracks, one track per object
\cite{bewley2016sort,zhang2022bytetrack,cao2023ocsort}. The tracker uses a box only if its score
passes a fixed limit, called a threshold. A two-stage tracker, such as ByteTrack, first matches the
high-score boxes to its tracks and then uses lower-score boxes in a second step, only to continue
tracks that are still unmatched \cite{zhang2022bytetrack}. A single-stage tracker, such as OC-SORT,
uses one threshold and ignores every box below it \cite{cao2023ocsort}. Table~\ref{tab:terms}
explains the main terms of this paper in plain words.

In real systems, the detector and the tracker are rarely built together. A detector is replaced by
a faster model, retrained for a new camera, or made smaller (quantized) for embedded hardware. The
tracker, however, keeps the thresholds that its authors set for another detector. Detector scores
are not calibrated across models and conditions \cite{kuppers2020detcalib}: the same object can get
a high score from one detector and a much lower score from another. A new detector therefore
usually gives scores on a different scale, and nothing in the system notices the change. The effect
can be large. In this paper, halving the scores of a published detector reduces two well-known
trackers to 0 higher order tracking accuracy (HOTA) \cite{luiten2021hota}, the main tracking metric
of this paper, because no detection passes their fixed thresholds (Sect.~\ref{sec:shift}).
Retuning every tracker for every detector needs labels from the new setting, that is, boxes and
identities drawn by people, and a working system usually does not have them.

Existing adaptive methods only partly address this problem. Trackers with adaptive thresholds
change their own thresholds from the detections of each frame, but each method is built into one
tracker \cite{ma2023adaptivebytetrack,shim2024adaptrack,act2026adaptive}. Score calibration methods
correct the detector scores, but they need labelled data \cite{kuppers2020detcalib}. Methods that
adapt a vision system while it runs learn or measure their settings for one model
\cite{han2022dynamic,jiang2018chameleon}. What is missing is a simple, training-free component that
works outside the tracker, adjusts itself from the detector's own scores, and leaves the tracker
alone when it already fits those scores.

This paper presents such a component: a self-calibrating control layer (\scl{}) that sits between
a detector and a tracker that are both kept fixed (frozen). We call the tracker that \scl{} serves
the \emph{host}, because \scl{} sits in front of it and works for it without changing its code. The
host gives \scl{} four numbers from its published settings, which we call the host contract
(Sect.~\ref{sec:contract}). The central idea is that the detector's own scores show whether its
output fits the host. \scl{} looks at the scores of recent frames and checks how clearly they
separate confident objects from clutter, that is, boxes that are not real objects. If the
separation is clear, the host keeps its own settings, and the layer only adds help that the host
lacks, such as a step that continues tracks. If it is not clear, \scl{} chooses which candidates
reach the host, gives them new scores, and sets the host's thresholds. What is new is therefore not
a detector or a tracker, but a layer outside the tracker that needs no labels, uses one fixed
configuration for many trackers, and stays out of the way when the tracker already fits the
detector.

\begin{table*}[t]
\caption{Main terms of this paper in plain words.}
\label{tab:terms}
\centering\footnotesize
\setlength{\tabcolsep}{4pt}
\renewcommand{\arraystretch}{1.15}
\begin{tabular}{@{}>{\raggedright\arraybackslash}p{32mm}>{\raggedright\arraybackslash}p{123mm}@{}}
\toprule
Term & Meaning in this paper\\
\midrule
Detector & A neural network that finds objects in one video frame and returns a box for each of them.\\
Candidate (detection) & One box returned by the detector, with a class label (for example, car) and a confidence score.\\
Confidence score & A number between 0 and 1 that says how sure the detector is that the box holds an object. Different detectors use different score scales.\\
Threshold & A fixed score limit. The tracker uses a candidate only if its score is at or above the limit.\\
Tracker & A program that links the candidates of consecutive frames into tracks, one track per object.\\
Track, identity (ID) & The boxes of one object over time, with one ID number that should never change.\\
Sequence & One video of a dataset.\\
Stream & The detector output of one video, frame after frame.\\
Host & The tracker that \scl{} serves. It runs after \scl{} with its own published code and settings, which are never changed.\\
Host contract & Four published settings of the host that \scl{} reads (Sect.~\ref{sec:contract}).\\
Matching and birth thresholds & The score a candidate needs to continue an existing track (matching, also called association), and the score it needs to start a new track (birth).\\
Single-stage, two-stage tracker & A single-stage tracker uses one threshold and ignores every candidate below it. A two-stage tracker also uses low-score candidates in a second step, only to continue existing tracks.\\
Frozen & Kept fixed: not retrained, not retuned, not changed in any way.\\
Labels (ground truth) & Boxes and identities drawn by people. They are used only to measure the results, never by \scl{}.\\
Development data & Data that was looked at while \scl{} was designed.\\
Held-out, predeclared & Data or trackers that were chosen in advance, were not used during the design, and were run once with the frozen \scl{}.\\
Logit & The score written on another scale, $\log\bigl(s/(1-s)\bigr)$, Eq.~\eqref{eq:logit}. It spreads out scores close to 0 and 1.\\
Otsu's method & A standard image-processing rule that picks the threshold that best splits a histogram into two groups \cite{otsu1979}.\\
HOTA, MOTA, IDF1 & Tracking metrics in percent; higher is better (Sect.~\ref{sec:metrics}). HOTA, the main metric, balances finding the objects and keeping their identities.\\
FP, FN, IDS & False positives (boxes where there is no object), false negatives (missed objects), and identity switches (an object's ID changes). They are counts; lower is better.\\
Paired comparison & The same tracker runs twice on the same sequences and detections, once alone and once with \scl{}.\\
95\% interval, significant & The range of likely values of a difference, found by resampling the test sequences 10,000 times (bootstrap). A difference is called significant when this range does not include zero.\\
\bottomrule
\end{tabular}
\end{table*}


The results are consistent across trackers, detectors, and data. With halved detector scores,
\scl{} restored two trackers from 0 to about 66 HOTA. It improved every single-stage tracker that
was tested, including one that was not seen during design, and it raised HOTA on every MOT17
sequence for both of them. With other detectors, it raised the multiple object tracking accuracy
(MOTA), a metric that counts errors, of two-stage trackers by up to 7.9 points on KITTI and by 24.7
points on held-out aerial data, that is, data that was not used during the design. It also cut the
identity switches, the cases where an object's ID changes, by 43\% and 47\% for the two-stage
trackers on KITTI and by 45\% and 56\% on the aerial data. On the same scores, it reached higher
HOTA than simple adaptive thresholds by 2.2 to 8.6 points, and it was ahead of the published
adaptive-threshold method ACT under halved scores (by 2.0 and 10.1 HOTA points) and on held-out
aerial data (by 1.8 HOTA points with YOLOv8n and 22.7 MOTA points with RT-DETR-L). Strong trackers
that already fit their detector kept their accuracy, and no statistically significant HOTA loss was
found in any tested combination of data, detector, and tracker.

The main contributions are:
\begin{enumerate}
\item \scl{}, a training-free and causal control layer that adjusts itself from the detector scores,
      keeps the tracker's settings when the scores are clean, and takes full control only when they
      are noisy, together with a host contract of four published settings that lets one fixed
      configuration work with single-stage and two-stage trackers without changing them
      (Sect.~\ref{sec:method}).
\item A development path from always-on banding (\aob{}), an earlier design that always
      intervened, to the final layer, in which each step is compared with the step it was built on
      (Sect.~\ref{sec:development}).
\item An evaluation of one frozen configuration on seven published trackers, four detection
      sources, and three datasets, with paired bootstrap intervals, a comparison with simple
      adaptive thresholds and with the published ACT method, and an analysis of each part of the
      layer (Sect.~\ref{sec:results}).
\item A latency measurement showing that the layer adds only about 3 ms per frame on an ordinary
      processor (central processing unit, CPU) (Sect.~\ref{sec:cost}).
\end{enumerate}

The rest of the paper is organized as follows. Section~\ref{sec:related} describes the datasets and
evaluation metrics and reviews related work. Section~\ref{sec:method} specifies \scl{} with all its
values, and Section~\ref{sec:development} explains how it was developed from an earlier design that
always intervened. The final configuration was then frozen before any external system was run.
Section~\ref{sec:results} gives the experimental settings and the results, and
Section~\ref{sec:discussion} a discussion. Section~\ref{sec:conclusion} concludes the paper.
